%%%%%%%%%%%%%%%%%%%%%%%%%%%%%%%%%%%%%%%%%%%%%%%%%%
%%%%    ~~~~ Abstract ~~~~
%%%%%%%%%%%%%%%%%%%%%%%%%%%%%%%%%%%%%%%%%%%%%%%%%%
\begin{abstract}
    Industrielle Qualitätskontrollsysteme und medizinische Bilddiagnostik lösen
    strukturell identische Aufgaben: die Lokalisation und Klassifikation von Anomalien
    in Bilddaten durch trainierte neuronale Netze. Der vorliegende Artikel untersucht,
    inwieweit die in der industriellen Bildanalyse etablierten Verfahren -- insbesondere
    Convolutional Neural Networks und Transfer Learning -- auf die medizinische Bildgebung
    übertragbar sind. Als empirischer Beleg dient das Modell CheXNet der Stanford
    University, das bei der Pneumonieerkennung auf Röntgenaufnahmen den durchschnittlichen
    F1-Score praktizierender Radiologen statistisch signifikant übertraf (0{,}435 vs.\
    0{,}387) \cite{chexnet2017}. Der Einsatz von KI als ergänzender zweiter Leser steigert
    die Sensitivität um den Faktor 1{,}12 bei gleichzeitiger Reduktion des Befundungsaufwands
    um rund 44~\% \cite{breastScreening2025}. Zur Quantifizierung des Nutzens werden
    standardisierte Metriken (Sensitivität, Spezifität, F1-Score, AUC-ROC) eingesetzt.
    Da medizinische Bilddaten besondere Schutzwürdigkeit genießen (Art.~9 DSGVO) und
    diagnostische KI-Systeme dem EU AI Act als Hochrisiko-KI unterliegen, werden
    datenschutzkonforme Systemarchitekturen -- insbesondere Federated Learning mit
    Differential Privacy sowie On-Premise-Infrastrukturen -- als Implementierungsoptionen
    analysiert \cite{federatedLearning2022, euAiAct2024}.
\end{abstract}
